\section{Computational Methods}
\label{sec:methods}

% REVISED: Concise methods for PRB Rapid Communications.
% Goal: sufficient detail for reproducibility without excessive length.

We perform density functional theory (DFT) calculations using
\textsc{Quantum ESPRESSO}~\cite{QE2020} with the PBEsol
exchange-correlation functional~\cite{PBEsol2008}, which yields more
accurate lattice parameters and bulk moduli under pressure than
PBE~\cite{Csonka2009}. We employ optimized norm-conserving Vanderbilt
(ONCV) pseudopotentials from the PseudoDojo library~\cite{PseudoDojo2018}
at stringent accuracy, as required for compatibility with the Wannier
interpolation step in EPW. The plane-wave kinetic energy cutoff is set
to $E_{\mathrm{cut}} = 80\text{--}100$~Ry depending on the element
combination, with convergence verified to $< 1$~mRy/atom.
Self-consistent field (SCF) calculations use a $\Gamma$-centered
$24 \times 24 \times 24$ Monkhorst-Pack~\cite{MonkhorstPack1976}
$k$-point grid and Methfessel-Paxton smearing of 0.02~Ry. Crystal
structures at each target pressure are obtained from variable-cell
relaxation with force and stress convergence thresholds of
$10^{-4}$~Ry/Bohr and 0.5~kbar, respectively.

Phonon frequencies and electron-phonon matrix elements are computed
within density functional perturbation theory
(DFPT)~\cite{Baroni2001} on a coarse $q$-point grid. We systematically
test convergence across $4 \times 4 \times 4$, $6 \times 6 \times 6$,
and $8 \times 8 \times 8$ grids, finding that $\lambda$ is converged
to within 5\% at $6 \times 6 \times 6$. The acoustic sum rule is
enforced via \texttt{asr=`crystal'} in \texttt{matdyn.x} to eliminate
spurious imaginary frequencies at $\Gamma$ arising from numerical noise.

Electron-phonon coupling properties are computed using the EPW
code~\cite{Ponce2016,Ponce2023}, which interpolates DFPT matrix
elements onto fine Brillouin zone meshes via maximally localized
Wannier functions~\cite{Marzari1997} constructed with
\textsc{Wannier90}~\cite{Wannier90}. The disentanglement window
spans $\pm 3$~eV around the Fermi level to capture all bands
contributing to $\lambda$. We employ fine interpolation grids of
$40^3$ for $k$-points and $20^3$ for $q$-points, verifying that
$\lambda$ changes by less than 2.5\% when increasing to $60^3$/$30^3$.
The Fermi surface thickness parameter is set to
\texttt{fsthick}~$= 0.4$~eV, with Gaussian broadening
$\sigma = 0.025$~eV for the Fermi-surface $\delta$-functions.

The superconducting critical temperature $T_c$ is determined by
solving the isotropic Eliashberg equations on the Matsubara
axis~\cite{Marsiglio2003}, as implemented in EPW. The Coulomb
pseudopotential $\mu^*$ is \emph{not} treated as an adjustable
parameter; instead, we report $T_c$ at fixed $\mu^* = 0.10$ and
$0.13$, bracketing the standard range for metallic
hydrides~\cite{Allen1975}. We verify consistency with the
Allen-Dynes modified McMillan formula~\cite{Allen1975} as a
cross-check; for $\lambda > 2$, the Allen-Dynes result systematically
underestimates the Eliashberg $T_c$ by 28--31\%, as expected in the
strong-coupling regime.

Anharmonic phonon renormalization is treated within the stochastic
self-consistent harmonic approximation
(SSCHA)~\cite{Monacelli2021,Errea2014} using the
\textsc{python-sscha} package with \textsc{Quantum ESPRESSO} as the
force engine. We employ a $2 \times 2 \times 2$ supercell (40~atoms)
with 100 stochastically displaced configurations per population
(200 for the 3~GPa quantum-stabilization analysis) at $T = 0$~K to
isolate quantum zero-point effects. The SSCHA minimization is
iterated over 20 populations until the free energy gradient converges
below $5 \times 10^{-8}$~Ry$^2$ and the free energy fluctuation
over the last three populations is below 0.05~meV/atom. The Kong-Liu
effective sample ratio is maintained above 0.5 throughout. The
anharmonic Eliashberg spectral function $\alpha^2F_{\mathrm{anh}}(\omega)$
is constructed by rotating the DFPT electron-phonon matrix elements
into the SSCHA eigenvector basis and evaluating at the renormalized
SSCHA frequencies~\cite{Errea2015,Belli2025}. The anharmonic coupling
constant decomposes as
$\lambda_{\mathrm{anh}} = R_{\mathrm{freq}} \times R_{\mathrm{rot}} \times \lambda_{\mathrm{harm}}$,
where $R_{\mathrm{freq}}$ captures the frequency hardening and
$R_{\mathrm{rot}}$ the eigenvector rotation effect.

Thermodynamic stability is assessed via the convex hull distance
$E_{\mathrm{hull}}$ (meV/atom) in the Cs--In--H ternary phase space
at each target pressure, computed from DFT enthalpies of all competing
elemental and binary phases using
\textsc{pymatgen}~\cite{Ong2013}. Dynamic stability is verified from
the absence of imaginary phonon frequencies across the full Brillouin
zone after SSCHA renormalization. The validity of the Migdal
approximation is confirmed by $\omega_{\log}/E_F = 0.009$--$0.013
\ll 0.1$ at all pressures studied.

%% CITATIONS NEEDED (for gpd-bibliographer):
%% - QE2020 — Giannozzi et al., J. Chem. Phys. 152, 154105 (2020), QE reference
%% - PBEsol2008 — Perdew et al., PRL 100, 136406 (2008), PBEsol functional
%% - Csonka2009 — Csonka et al., PRB 79, 155107 (2009), PBEsol benchmarks for solids
%% - PseudoDojo2018 — van Setten et al., Comp. Phys. Commun. 226, 39 (2018)
%% - MonkhorstPack1976 — Monkhorst and Pack, PRB 13, 5188 (1976)
%% - Baroni2001 — Baroni et al., Rev. Mod. Phys. 73, 515 (2001), DFPT review
%% - Ponce2016 — Ponce et al., Comp. Phys. Commun. 209, 116 (2016), EPW code
%% - Ponce2023 — Ponce et al., npj Comput. Mater. 9, 170 (2023), EPW update
%% - Marzari1997 — Marzari and Vanderbilt, PRB 56, 12847 (1997), MLWFs
%% - Wannier90 — Pizzi et al., J. Phys.: Condens. Matter 32, 165902 (2020)
%% - Marsiglio2003 — Marsiglio and Carbotte, in The Physics of Superconductors (2003)
%% - Allen1975 — Allen and Dynes, PRB 12, 905 (1975)
%% - Monacelli2021 — Monacelli et al., J. Phys.: Condens. Matter 33, 363001 (2021)
%% - Errea2014 — Errea et al., PRB 89, 064302 (2014), SSCHA original
%% - Errea2015 — Errea et al., PRL 114, 157004 (2015), SSCHA for H3S
%% - Belli2025 — Belli et al., arXiv:2507.03383 (2025), YH6 SSCHA benchmarks
%% - Ong2013 — Ong et al., Comp. Mater. Sci. 68, 314 (2013), pymatgen
